%%%%%%%%%%  scan idx 335  |  folio 322  %%%%%%%%%%
\SGAPrintPage{322}%
whose commutativity is verified, by 1.10, on the elements \(1+Xt,1+Yt\). For the first,
\[
\begin{aligned}
\lambda_u(1+Xt)\circ\lambda_u(1+Yt)
&=(\underline{1}+\{1+Xt\}u)\circ(\underline{1}+\{1+Yt\}u)\\
&=\underline{1}+\{(1+Xt)\circ(1+Yt)\}u\\
&=\underline{1}+\{1+XYt\}u\\
&=\lambda_u((1+Xt)\circ(1+Yt)).
\end{aligned}
\]
For the second:
\[
  \lambda_v(\lambda_u(1+Xt))
  =
  \underline{1}+\{\underline{1}+\{1+Xt\}u\}v
  =
  \Ghat^v(\lambda_u)(\lambda_v(1+Xt)).
\]

\item The geometric examples of 2.2 will be treated later.
\end{enumerate}

\sgaStableTarget{sga6.e05.para.2-6}2.6. We can generalize Example 2.5.a) as follows: let \(K\) be a commutative ring with unit, such that:
\[
\begin{array}{ll}
\text{i)}& K\text{ is torsion-free},\\
\text{ii)}& \forall x\in K,\ \forall n\in\mathbb N,\quad \text{we have}\ \binom{x}{n}\in K
\end{array}
\]
We then put, for every \(x\in K\), and every \(n\in\mathbb N\), \(\lambda^n(x)=\binom{x}{n}\), that is:
\sgaStableTarget{sga6.e05.eq.2-6-1}\begin{equation}\tag{2.6.1}
  \lambda_t(x)=\sum_{n=0}^{\infty}\binom{x}{n}t^n
  =(1+t)^x=\exp(x\Log(1+t)),
\end{equation}
where the last identity is considered in \((K\otimes_{\mathbb Z}\mathbb Q)[[t]]\). By the multiplicativity of the exponential, this defines a pre-\(\lambda\)-ring structure on \(K\).

To see that it is a \(\lambda\)-ring structure, we must prove relations 2.4.2, and it suffices to prove them in \(K\otimes_{\mathbb Z}\mathbb Q\), by 2.6.i). Functoriality then reduces us to the case \(K=\mathbb Q[X,Y]\) and the elements \(X,Y\). In view of the definition of the \(\lambda^i\), the two members of the
%%%%%%%%%%  scan idx 336  |  folio 323  %%%%%%%%%%
\SGAPrintPage{323}%
equalities 2.4.2 are polynomials in \(X,Y\), with rational coefficients, and to prove their equality it therefore suffices to show that they take the same value at all integer values of \(X,Y\). But the relations obtained by substituting integer values for \(X,Y\) are those expressing that the pre-\(\lambda\)-ring \(Z\) (2.5.a)) is a \(\lambda\)-ring; they are therefore satisfied.

\sgaStableTarget{sga6.e05.def.2-7}\textbf{Definition 2.7.} \uline{By a} binomial ring \uline{we mean a ring satisfying conditions} 2.6.i) \uline{and} ii), \uline{endowed with the} \(\lambda\)-\uline{ring structure defined by} 2.6.1.

\sgaStableTarget{sga6.e05.ex.2-8}\textbf{Example 2.8.} Let \(X\) be a site, \(K\) a binomial ring, and \(\underline{K}\) the sheaf associated to the constant presheaf with value \(K\) on \(X\); then \(H^0(X,\underline{K})\) is a binomial ring.

Indeed, conditions 2.6 are preserved under passage to an inductive limit.

In particular, \(H^0(X,\underline{\mathbb Z})\) is a binomial ring.

\sgaStableTarget{sga6.e05.para.2-9}2.9. Let \(K\) be a binomial ring, and \(A\) a \(K\)-algebra. For every series \(\varphi\in\Ghat(A)\), and every \(a\in K\), we can define a series \(\varphi^a\) in the usual way: we write \(\varphi=1+\psi\), and put:
\sgaStableTarget{sga6.e05.eq.2-9-1}\begin{equation}\tag{2.9.1}
  \varphi^a=(1+\psi)^a
  =
  \sum_{n=0}^{\infty}\binom{a}{n}\psi^n
  \left(=\exp(a\Log\varphi)\right),
\end{equation}
the last equality being writable only modulo torsion. These formulas are obtained by substitution from the analogous formulas for \(\varphi=1+t\); the following relations, the classical properties of the exponential, therefore remain valid for \(a,b\in K\), \(\varphi,\psi\in\Ghat(A)\):
\sgaStableTarget{sga6.e05.eq.2-9-2}\begin{equation}\tag{2.9.2}
  \varphi^a\varphi^b=\varphi^{a+b};\qquad
  \varphi^{ab}=(\varphi^a)^b;\qquad
  \varphi^a\psi^a=(\varphi\psi)^a .
\end{equation}

% Native English continuation macros retained at their source boundary.
\providecommand{\Ghat}{\widehat G}
\providecommand{\Gg}{\widehat G_g}
\providecommand{\Gcirc}{\widehat G_{\circ}}
\providecommand{\Gnabla}{\widehat G_{\nabla}}
\providecommand{\Grhat}{\widehat{\operatorname{Gr}}}
\providecommand{\Log}{\operatorname{Log}}
\providecommand{\Fil}{\operatorname{Fil}}
\providecommand{\eps}{\varepsilon}
\providecommand{\id}{\operatorname{id}}
\providecommand{\Sym}{\operatorname{Sym}}
\providecommand{\Mon}{\operatorname{Mon}}

% Exposé V continuation macros
\providecommand{\Ghat}{\widehat G}
\providecommand{\Gg}{\widehat G_g}
\providecommand{\Gcirc}{\widehat G_{\circ}}
\providecommand{\Gnabla}{\widehat G_{\nabla}}
\providecommand{\Grhat}{\widehat{\operatorname{Gr}}}
\providecommand{\Log}{\operatorname{Log}}
\providecommand{\Fil}{\operatorname{Fil}}
\providecommand{\eps}{\varepsilon}
\renewcommand{\id}{\operatorname{id}}
\providecommand{\Sym}{\operatorname{Sym}}
\providecommand{\Mon}{\operatorname{Mon}}

%%%%%%%%%%  scan idx 337  |  folio 324  %%%%%%%%%%
\SGAPrintPage{324}%
\sgaStableTarget{sga6.e05.continuation-of-2-9}We still suppose that \(a\in K\). Multiplication by \(\lambda_t(a)\) under the law \(\circ\) defines a morphism of functors: \(\Ghat_K\to\Ghat_K\), characterized by its value on the element \(1+Xt\) of \(\Ghat(K[X])\). We therefore have:
\[
  \lambda_t(a)\circ(1+Xt)=\lambda_{Xt}(a)\qquad\text{by (2.3.2), hence:}
\]
\sgaStableTarget{sga6.e05.eq.2-9-3}\begin{equation}\tag{2.9.3}
  \lambda_t(a)\circ(1+Xt)
  =(1+Xt)^a
  =\exp(a\Log(1+Xt))
  =\sum_{n=0}^{\infty}\binom{a}{n}X^nt^n .
\end{equation}
We deduce:
\sgaStableTarget{sga6.e05.eq.2-9-4}\begin{equation}\tag{2.9.4}
  \lambda_t(a)\circ\varphi=\varphi^a
  \quad(=\exp(a\Log\varphi)) ;
\end{equation}
Indeed, the two members are additive functors equal on \(1+Xt\). From (2.9.3) and 1.15, it follows that the \(n\)-th universal polynomial defining multiplication by \(\lambda_t(a)\) is isobaric of weight \(n\); it will be denoted by \(R_{n,a}\).

Finally, note that if \(A\) is a \(K\)-\(\lambda\)-algebra, i.e. a \(\lambda\)-ring endowed with a \(\lambda\)-homomorphism: \(K\to A\), and if \(b\in A\), we have:
\sgaStableTarget{sga6.e05.eq.2-9-5}\begin{equation}\tag{2.9.5}
  \lambda_t(ab)=(\lambda_t(b))^a .
\end{equation}
This follows from \(\lambda_t(ab)=\lambda_t(a)\circ\lambda_t(b)\), and from (2.9.4).

\sgaStableTarget{sga6.e05.sec.3}\section*{3. The operations \(\gamma\)}

We still work in the category \(\mathcal C_0\).

\sgaStableTarget{sga6.e05.para.3-1}3.1. Let \(K\) be a \(\lambda\)-ring. From the operations \(\lambda^i\) of \(K\), we shall define new operations, denoted by \(\gamma^i\). They are intended to introduce a notion of Chern classes under conditions that will be specified later. For the moment, let us give a ``justification'' of the following definitions: if \(\underline{E}\) is a vector bundle of rank \(d\), the \(d\)-th Chern class of \(\underline{E}\) will have to be (whatever theory of Chern classes is considered) the ``alternating sum of the exterior powers
%%%%%%%%%%  scan idx 338  |  folio 325  %%%%%%%%%%
\SGAPrintPage{325}%
of \(\underline{E}\)'';
to define the \(n\)-th Chern class, we shall proceed analogously, bringing \(\underline{E}\) to rank \(n\) by considering \((\cl(\underline{E})-d)+n\), and we shall take the alternating sum of the exterior powers of the element thus obtained; the following proposition is to be ``interpreted'' by setting \(x=\cl(\underline{E})-d\):

\sgaStableTarget{sga6.e05.prop.3-2}\textbf{Proposition 3.2.} \uline{Let} \(K\) \uline{be a} \(\lambda\)-\uline{ring}.
\begin{enumerate}[label=\roman*)]
\item \(\forall x\in K,\ \forall n\in\mathbb N,\quad \lambda^n(x+n-1)=(-1)^n\sum_{i=0}^n(-1)^i\lambda^i(x+n).\)
\item \uline{Putting} \(\gamma^n(x)=\lambda^n(x+n-1)\), \uline{and} \(\gamma_t(x)=\sum_{n\ge0}\gamma^n(x)t^n\), \uline{we have}:
\[
  \gamma_t(x)=\lambda_{t/(1-t)}(x),\qquad\text{\uline{hence}}
\]
\[
  \lambda_t(x)=\gamma_{t/(1+t)}(x).
\]
\end{enumerate}

We therefore see that knowing the operations \(\gamma\) is equivalent to knowing the operations \(\lambda\).

\uline{Proof}.

i) We have
\[
  \lambda^n(x+n-1)=\sum_{i=0}^{n}\lambda^i(x+n)\lambda^{n-i}(-1).
\]
From \(\lambda_t(1)=1+t\), we obtain: \(\lambda_t(-1)=1/(1+t)\), hence: \(\lambda^k(-1)=(-1)^k\), which gives the required relation.

ii) \(\lambda_{t/(1-t)}(x)=\sum_{n\ge0}\lambda^n(x).t^n(1-t)^{-n}.\)
We expand \(t^n(1-t)^{-n}=t^n\sum_{k\ge0}\binom{-n}{k}(-1)^kt^k=\sum_{p\ge n}\binom{p-1}{p-n}t^p\). \(\lambda_{t/(1-t)}(x)=\sum_{n\ge0}\bigl[\lambda^n(x).\sum_{p\ge n}\binom{p-1}{p-n}t^p\bigr]\); the coefficient of \(t^p\) is therefore:
\[
  \sum_{n=0}\lambda^n(x).\binom{p-1}{p-n}
  =\sum_{n=0}\lambda^n(x)\lambda^{p-n}(p-1)
  =\lambda^p(x+p-1).
\]

%%%%%%%%%%  scan idx 339  |  folio 326  %%%%%%%%%%
\SGAPrintPage{326}%
\sgaStableTarget{sga6.e05.para.3-3}3.3.
Let \(K\), \(K'\) be two \(\lambda\)-rings, and \(f:K\to K'\) a ring homomorphism. By Definition 3.2.ii), \(f\) is a \(\lambda\)-homomorphism if and only if, for every \(i\in\mathbb N\), and every \(x\in K\), \(f(\gamma^i(x))=\gamma^i(f(x))\), or equivalently \(\gamma_t(f(x))=\Gg(f)\circ\gamma_t(x)\).

\sgaStableTarget{sga6.e05.para.3-4}3.4. From relation 3.2.ii), we deduce the relations:
\sgaStableTarget{sga6.e05.eq.3-4-1}\begin{equation}\tag{3.4.1}
  \gamma^0(x)=1,
  \qquad \gamma^1(x)=x,
  \qquad \gamma^n(x+y)=\sum_{p+q=n}\gamma^p(x)\gamma^q(y).
\end{equation}
They show that the homomorphism \(\gamma_t\) defines on \(K\) a new pre-\(\lambda\)-ring structure; we shall see that in general it is not a \(\lambda\)-ring structure. Denoting by \(\rho_t\) the substitution in the ring of formal series in \(t\) with coefficients in \(K\) defined by \(\rho_t(t)=t/(1-t)\), we have:
\[
  \gamma_t=\rho_t\circ\lambda_t.
\]
The automorphism \(\rho_t\) is functorial in \(K\), and defines an automorphism of the functor \(\Ghat\) considered as a functor to abelian groups; with this automorphism, by transport of structure, we can define on \(\Ghat\) a new structure as a functor to \(\lambda\)-rings, such that for every \(\lambda\)-ring \(K\), the homomorphism \(\gamma_t:K\to\Ghat(K)\) is a \(\lambda\)-homomorphism; we shall make explicit the multiplication and the operations \(\lambda^i\) of this second structure.

\sgaStableTarget{sga6.e05.para.3-5}3.5. The multiplication is defined by the homomorphism of functors:
\[
  \Ghat\times\Ghat\xrightarrow{(\rho^{-1},\rho^{-1})}
  \Ghat\times\Ghat\xrightarrow{\circ}\Ghat\xrightarrow{\rho}\Ghat ,
\]
where the sign \(\circ\) denotes the multiplication defined in 2.3. This new multiplication will be denoted by \(v\). As usual, it is defined by its value on \((1+Xt,1+Yt)\); we therefore have:
